\MathBookSourcePage{169}
\subsection{网格依赖范数}
\label{sec:ch06-mesh-dependent-norms}
\setcounter{thmcount}{0}
\setcounter{equation}{0}

\MBSourceItem{1}
\begin{Definition}[网格依赖内积]
\label{def:ch06-mesh-dependent-inner-product}
$V_k$ 上的网格依赖内积 $(\cdot,\cdot)_k$ 定义为
\MBSourceItem{2}
\begin{equation}
\MathBookFormulaID{formula-ch06-mesh-dependent-inner-product}
\label{eq:ch06-mesh-dependent-inner-product}
(v,w)_k:=h_k^2\sum_{i=1}^{n_k}v(p_i)w(p_i),
\end{equation}
其中 $\{p_i\}_{i=1}^{n_k}$ 是 $\mathcal T_k$ 的内部顶点集合.
\end{Definition}

算子 $A_k:V_k\longrightarrow V_k$ 定义为
\MBSourceItem{3}
\begin{equation}
\MathBookFormulaID{formula-ch06-discrete-operator-definition}
\label{eq:ch06-discrete-operator-definition}
(A_kv,w)_k=a(v,w)\qquad\forall v,w\in V_k.
\end{equation}
以算子 $A_k$ 表示时, 离散方程~\eqref{eq:ch06-discrete-problem} 可写为
\MBSourceItem{4}
\begin{equation}
\MathBookFormulaID{formula-ch06-discrete-operator-equation}
\label{eq:ch06-discrete-operator-equation}
A_ku_k=f_k,
\end{equation}
其中 $f_k\in V_k$ 满足
\MBSourceItem{5}
\begin{equation}
\MathBookFormulaID{formula-ch06-discrete-load-vector}
\label{eq:ch06-discrete-load-vector}
(f_k,v)_k=(f,v)\qquad\forall v\in V_k.
\end{equation}
由于 $A_k$ 关于 $(\cdot,\cdot)_k$ 对称正定, 可按通常方式定义一族网格依赖范数 $\|\cdot\|_{s,k}$:
\MBSourceItem{6}
\begin{equation}
\MathBookFormulaID{formula-ch06-mesh-dependent-norm-scale}
\label{eq:ch06-mesh-dependent-norm-scale}
\|v\|_{s,k}:=\sqrt{(A_k^sv,v)_k},
\end{equation}
其中 $A_k^s$ 表示对称正定算子 $A_k$ 的 $s$ 次幂(Halmos).

网格依赖范数 $\|\cdot\|_{s,k}$ 将在收敛分析中发挥重要作用. 注意能量范数 $\|\cdot\|_E=\sqrt{a(\cdot,\cdot)}$ 与 $V_k$ 上的 $\|\cdot\|_{1,k}$ 范数一致. 类似地, $\|\cdot\|_{0,k}$ 是与网格依赖内积~\eqref{eq:ch06-mesh-dependent-inner-product} 相伴的范数. 下一个引理说明 $\|\cdot\|_{0,k}$ 与 $L^2$ 范数等价.

\MBSourceItem{7}
\begin{Lemma}
\label{lem:ch06-l2-mesh-norm-equivalence}
存在正常数 $C_1,C_2$, 使得
\[
\MathBookFormulaID{formula-ch06-l2-mesh-norm-equivalence}
C_1\|v\|_{L^2(\Omega)}\leq\|v\|_{0,k}\leq C_2\|v\|_{L^2(\Omega)}
\qquad\forall v\in V_k.
\]
\end{Lemma}
\begin{Proof}
利用三角形上二次多项式的求积公式,
\[
\MathBookFormulaID{formula-ch06-edge-midpoint-quadrature}
\begin{aligned}
\|v\|_{L^2(\Omega)}^2
&=\sum_{T\in\mathcal T_k}\int_Tv^2\,dx\\
&=\sum_{T\in\mathcal T_k}\frac{|T|}{3}
\left(\sum_{i=1}^3v^2(m_i)\right)
\qquad(\text{参见习题~\ref{exr:ch06-12}}),
\end{aligned}
\]
\MathBookSourcePage{170}
其中 $|T|$ 表示 $T$ 的面积. 由定义,
\[
\MathBookFormulaID{formula-ch06-mesh-norm-vertex-sum}
\|v\|_{0,k}^2=(v,v)_k=h_k^2\sum_{i=1}^{n_k}v^2(p_i).
\]
$\{\mathcal T_k\}$ 的拟一致性蕴含 $|T|$ 与 $h_k^2$ 等价. 因为 $v$ 是线性的,
\[
\MathBookFormulaID{formula-ch06-midpoint-vertex-transform}
\begin{pmatrix}v(m_1)\\v(m_2)\\v(m_3)\end{pmatrix}
=
\begin{pmatrix}
0&\frac12&\frac12\\
\frac12&0&\frac12\\
\frac12&\frac12&0
\end{pmatrix}
\begin{pmatrix}v(p_1)\\v(p_2)\\v(p_3)\end{pmatrix},
\]
其中 $p_1,p_2,p_3$ 是 $T$ 的顶点, $m_i$ 是 $p_i$ 对边的中点, $i=1,2,3$. 引理得证.
\end{Proof}

下面估计 $A_k$ 的谱半径 $\Lambda(A_k)$.

\MBSourceItem{8}
\begin{Lemma}
\label{lem:ch06-spectral-radius-bound}
$\Lambda(A_k)\leq C h_k^{-2}$.
\end{Lemma}
\begin{Proof}
设 $\lambda$ 是 $A_k$ 的特征值, $\phi$ 是相应特征向量. 则
\[
\MathBookFormulaID{formula-ch06-spectral-eigenvalue-identities}
\begin{aligned}
a(\phi,\phi)&=(A_k\phi,\phi)_k\\
&=\lambda(\phi,\phi)_k\\
&=\lambda\|\phi\|_{0,k}^2.
\end{aligned}
\]
因此
\[
\MathBookFormulaID{formula-ch06-eigenvalue-rayleigh-quotient}
\lambda=\frac{a(\phi,\phi)}{\|\phi\|_{0,k}^2}.
\]
由于 $\|\phi\|_{0,k}^2$ 与 $\|\phi\|_{L^2(\Omega)}^2$ 等价, 由引理~\ref{lem:ch04-local-inverse-estimate} 可得此引理:
\MBSourceItem{9}
\begin{equation}
\MathBookFormulaID{formula-ch06-finite-element-inverse-estimate}
\label{eq:ch06-finite-element-inverse-estimate}
\|v\|_E\leq C h_k^{-1}\|v\|_{L^2(\Omega)}
\qquad\forall v\in V_k.
\end{equation}
\end{Proof}

最后给出一个有用的引理, 其证明留作习题(参见习题~\ref{exr:ch06-01}).

\MBSourceItem{10}
\begin{Lemma}[广义 Cauchy--Schwarz 不等式]
\label{lem:ch06-generalized-cauchy-schwarz}
对任意 $t\in\mathbb R$,
\[
\MathBookFormulaID{formula-ch06-generalized-cauchy-schwarz}
|a(v,w)|\leq\|v\|_{1+t,k}\|w\|_{1-t,k}
\qquad\forall v,w\in V_k.
\]
\end{Lemma}
